\chapter{Colorful $k$-median}
\label{cap:colkmed}

This section adapts the $7.081(1+\eps)$-approximation of Krishnaswamy, Li, and Sandeep~\cite{KLS2018} for the $k$-median with outliers problem, presented in Chapter~\ref{cap:kmedout}, to the \emph{colorful $k$-median} problem with two colors. We follow the notation of Chapter~\ref{cap:colkcenter} and split the client set into a set $R$ of \emph{red} clients and a set $B$ of \emph{blue} clients, requiring at least $\rho$ red clients and at least $\beta$ blue clients to be covered.

\begin{problem}[Colorful $k$-median]
Given a metric space $(F \cup D, d)$, where $F$ is a set of facility locations, $D = R \cup B$ is a set of clients partitioned into a set $R$ of red clients and a set $B$ of blue clients, and $d : (D \cup F) \times (D \cup F) \to \mathbb{R}_+$ is a metric, an integer $k$ representing the number of facility locations that can be opened, and integers $\rho$ and $\beta$ representing the coverage requirements for the red and blue clients respectively, the goal is to find a set $S \subseteq F$ and a set $C \subseteq D$ such that $|S| \le k$, $|C \cap R| \geq \rho$, and $|C \cap B| \geq \beta$, and $S$ and $C$ minimize $\sum_{j \in C} d(j,S)$.
We assume there are no colocated facilities, but clients may be colocated with each other or with a facility.
\end{problem}

The algorithms follow the same process as the one of the Chapter~\ref{cap:kmedout}. Then, we define the notion of \emph{colorful extended instances} for colorful $k$-median instances.

\begin{definition}
A \emph{colorful extended instance} of the colorful $k$-median problem is denoted by a tuple $(F, R, B, d, k, \rho, \beta, S_0)$, where $(F,R,B,d,k,\rho,\beta)$ is an instance of the colorful $k$-median problem and $S_0 \subseteq F$ is a set of pre-opened facilities. A feasible solution for a colorful extended instance is a pair $(S,C)$, where $S \subseteq F$ is a set of facilities such that $|S| \leq k$ and $S_0 \subseteq S$, and $C \subseteq R \cup B$ is a set of clients such that $|C \cap R| \geq \rho$ and $|C \cap B| \geq \beta$.
\end{definition}

\subsection{The natural LP and integrality gaps}

The natural LP relaxation of the colorful $k$-median problem is obtained from~\ref{LP-basic} by replacing the single coverage constraint with two constraints, one for each color:
\begin{equation}\tag{$\text{LP}_{\text{basic}}'$} \label{LP-basic-col}
\begin{aligned}
\text{Minimize} \quad &\sum_{i \in F} \sum_{j \in D} x_{ij} \,d(i,j) \\
\text{subject to} \quad &\sum_{i \in F} y_i \le k, \\
&x_{ij} \le y_i, && \forall i \in F, \ j \in D,\\
&\sum_{i \in F} x_{ij} \le 1, && \forall j \in D,\\
&\sum_{j \in R} \sum_{i \in F} x_{ij} \ge \rho, \\
&\sum_{j \in B} \sum_{i \in F} x_{ij} \ge \beta, \\
&x \in[0,1]^{F\times D}, \\
&y \in [0,1]^F.
\end{aligned}
\end{equation}

The two gap examples of Chapter~\ref{cap:kmedout} (Figure~\ref{fig:gap_examples}) do not use the coverage constraint in any way that depends on colors: both instances can be reproduced with a single color without changing the argument. Hence the same unbounded integrality gap arises here, for the same two reasons (the cost of increasing a fractional facility to $1$, and the cost of decreasing a fractional facility to $0$), and the same three preprocessing ideas of Chapter~\ref{cap:kmedout} are needed: bounding the star cost, bounding the backup cost, and sparsification.

\subsection{Preprocessing the instance}

Since the preprocessing ideas only bound connection costs, and never reference the coverage requirement or the color of a client, the notion of sparse instances and the preprocessing machinery of Chapter~\ref{cap:kmedout} extend to the colorful setting essentially unchanged; only the bookkeeping of how many red and blue clients are covered needs to be carried along. We restate the relevant definitions and theorems below, renaming the sparsity parameter from $\rho$ to $\sigma$ to avoid clashing with the red-coverage requirement $\rho$.

\begin{definition}[Sparse Colorful Extended Instances]
    \label{def:sparse_extended_instances_col}
Given a colorful extended instance $I' = (F,R,B,d,k,\rho,\beta,S_0)$, parameters $\sigma, \delta \in (0, 1/2)$, and $U \ge 0$. Let $(S^*, C^*)$ be a feasible solution to $I'$ with cost at most $U$. We say that $I'$ is $(\sigma, \delta, U)$-sparse w.r.t.\ $(S^*, C^*)$ if:
\begin{enumerate}[label=\thedefinition\alph*)]
    \item \label{def:sparse_col_a} for every $i \in S^* \setminus S_0$, one has $\sum_{j \in C_i^*} d(j, i) \leq \sigma U$;
    \item \label{def:sparse_col_b} for every $p \in F \cup D$, one has $|\mathcal{B}(p, \delta \cdot d(p, S^*)) \cap C^*| \cdot d(p, S^*) \leq \sigma U$.
\end{enumerate}
\end{definition}

This is the same notion as Definition~\ref{def:sparse_extended_instances}, with $D = R \cup B$ and $\rho$ renamed to $\sigma$; in particular $C_i^*$ still denotes, for $i \in S^*$, the set of clients of $C^*$ having $i$ as their closest facility, regardless of color.

\begin{theorem}[Adapted from Theorem~\ref{thm:sparsification}]
\label{thm:sparsification_col}
Suppose we are given an instance $I = (F, R, B, d, k, \rho, \beta)$, parameters $\sigma, \delta \in (0, 1/2)$, and an upper bound $U$ on the cost of an optimal solution $(S^*, C^*)$ to $I$ (which is not given to us). There is an $n^{O(1/\sigma)}$-time algorithm that outputs $n^{O(1/\sigma)}$ colorful extended instances such that one of them, say $I'$, has the form $(F, D' \subseteq D, d, k, \rho' = |C^* \cap R \cap D'|, \beta' = |C^* \cap B \cap D'|, S_0 \subseteq S^*)$ and satisfies:
\begin{enumerate}[label=\thedefinition\alph*)]
    \item \label{thm:sparsification_col_a} $I'$ is $(\sigma, \delta, U)$-sparse w.r.t.\ the solution $(S^*, C^* \cap D')$;
    \item \label{thm:sparsification_col_b} $\frac{1-\delta}{1+\delta} \sum_{j \in C^* \setminus D'} d(j, S_0) + \sum_{j \in C^* \cap D'} d(j, S^*) \leq U$.
\end{enumerate}
\end{theorem}

\begin{proof}
As in the proof of Theorem~\ref{thm:sparsification}, first assume that we know an optimal solution $(S^*,C^*)$ of $I$, and let $(\kappa^*,c^*)$ be the nearest-facility-vector-pair for $S^*$. We initialize $D' \coloneqq R \cup B$ and $S_0 \coloneqq \emptyset$, and run the same two-phase construction as in the proof of Theorem~\ref{thm:sparsification}: first, we add to $S_0$ every facility $i \in S^*$ with $\sum_{j \in C_i^*} d(j,i) > \sigma U$, of which there are at most $1/\sigma$; then, while there exists $p \in F \cup D$ with $|\mathcal{B}(p, \delta \cdot d(p, S^*)) \cap C^* \cap D'| \cdot d(p, S^*) > \sigma U$, we add the facility of $S^*$ nearest to $p$ to $S_0$ and remove $\mathcal{B}(p, \delta \cdot d(p, S^*))$ from $D'$.

This construction never inspects the color of a client, either when deciding to add a facility to $S_0$ or when deciding to remove a ball of clients from $D'$; it only compares connection costs against $\sigma U$. Hence the same argument as in the proof of Theorem~\ref{thm:sparsification} shows: the first phase adds at most $1/\sigma$ facilities to $S_0$; the second phase terminates after $O(1/\sigma)$ iterations, since each iteration removes clients whose combined cost exceeds $(1-\delta)\sigma U$ out of a total budget of $U$; and, at termination, Properties~\ref{thm:sparsification_col_a} and~\ref{thm:sparsification_col_b} hold for the instance $I' = (F, D', d, k, \rho' = |C^* \cap R \cap D'|, \beta' = |C^* \cap B \cap D'|, S_0)$ with respect to $(S^*, C^* \cap D')$, by the same triangle-inequality computation as before.

The only new bookkeeping is that, at the end, $I'$ additionally records how many of the surviving clients in $D'$ are red and how many are blue, namely $\rho'$ and $\beta'$.

Since $(S^*,C^*)$ is not known in practice, we cannot run this procedure directly. However, as before, the set $S_0$ produced has size $O(1/\sigma)$, and $D'$ is obtained from $D$ by removing $O(1/\sigma)$ balls, each determined by a center and one of $O(n^2)$ possible radii; there are also at most $n$ possible values for $\rho'$ and $n$ for $\beta'$. Thus there are at most $n^{O(1/\sigma)}$ possible instances $(F,D',d,k,\rho',\beta',S_0)$, and we simply output all of them; one of them is guaranteed to be the instance produced by running the procedure on the actual optimal solution.
\end{proof}

The construction of the vector of radius bounds $(R_j)_{j \in D}$ (Theorem~\ref{thm:radius_existence}) also does not reference the coverage requirement or the color of the clients at any point, so it applies verbatim to sparse colorful extended instances, with $\rho$ replaced by $\sigma$ throughout the statement and the proof; we do not restate it here.

\subsection{The colorful iterative rounding framework}

\subsubsection{The strengthened LP and eliminating the $x$ variables}

The strengthened LP~\eqref{LP-strong} is adapted analogously, by strengthening~\eqref{LP-basic-col} instead of~\eqref{LP-basic}:
\begin{align} \tag{$\text{LP}_{\text{strong}}'$} \label{LP-strong-col}
\text{Minimize} \quad &\sum_{i \in F} \sum_{j \in D} x_{ij} d(i,j) \\
\text{subject to} \quad &\text{constraints of } \eqref{LP-basic-col}, \nonumber \\
& y_i = 1, \quad &&\forall i \in S_0, \nonumber \\
& x_{ij} = 0, \quad &&\forall i \in F, \ j \in D \text{ s.t. } d(i,j) > R_j, \nonumber \\
& \sum_{j \in D} x_{ij} d(i,j) \leq \sigma \tilde{U} y_i, \quad &&\forall i \in F \setminus S_0, \nonumber \\
& x_{ij} = 0, \quad &&\forall i \in F \setminus S_0, \ j \in D \text{ s.t. } d(i,j) > \sigma \tilde{U}. \nonumber
\end{align}
where, as before, $\tilde{U} \coloneqq (1+\delta/2)U$. Since none of these constraints reference the coverage requirement, Lemma~\ref{lem:strong_lp_cost} (the cost of the optimum of the strengthened LP is at most $\tilde{U}'$) holds verbatim for~\eqref{LP-strong-col}.

Similarly, the elimination of the $x$ variables via collocated copies of facilities does not depend on colors, except for the bookkeeping of the coverage property:

\begin{lemma}[Adapted from Lemma~\ref{lem:eliminate_x}]
\label{lem:eliminate_x_col}
By adding collocated facilities to $F$, we can efficiently obtain a vector $y^* \in [0,1]^F$ and a set of allowed facilities $F_j \subseteq \mathcal{B}(j, R_j) \cap F$ for every client $j \in D$, satisfying the following properties:
\begin{itemize}
    \item[(a)] For each client $j \in D$, we have $y^*(F_j) \leq 1$;
    \item[(b)] The total fractionally open facilities satisfies $y^*(F) \leq k$;
    \item[(c)] The total fractional coverage of each color satisfies $\sum_{j \in R} y^*(F_j) \geq \rho$ and $\sum_{j \in B} y^*(F_j) \geq \beta$;
    \item[(d)] The bounded connection cost satisfies $\sum_{j \in D} \sum_{i \in F_j} d(i, j) y^*_i \leq \tilde{U}'$;
    \item[(e)] For each pre-opened facility $i \in S_0$, we have $\sum_{i' \text{ collocated with } i} y^*_{i'} = 1$;
    \item[(f)] For every facility $i$ that is not collocated with any facility in $S_0$, the star-cost satisfies $\sum_{j \in D : i \in F_j} d(i, j) y^*_i \leq 2\sigma\tilde{U}$.
\end{itemize}
\end{lemma}

The construction (splitting fractional facilities into collocated copies, ordered by star cost) and the proof of properties (a), (b), (d), (e), and (f) are identical to those of Lemma~\ref{lem:eliminate_x}, since they never reference the color of a client. Property (c) is inherited directly from the two coverage constraints of~\eqref{LP-strong-col}, one per color, in place of the single coverage constraint of~\eqref{LP-strong}.

\subsubsection{Random discretization of distances}

The randomized discretization of distances (Section~\ref{random_discretization_of_distances}) is also unaffected by colors: it only depends on the metric $d$. We use the same discretized distance $d'$ and the same constant $\tau \approx 2.3603$, and Lemma~\ref{lem:random_dist} applies without any change.

\subsubsection{Auxiliary LP and iterative rounding}

As in Section~\ref{kmedout}, the algorithm maintains a partition of $D = R \cup B$ into $D_\text{full}$, the set of clients which need to be fully connected, and $D_\text{part}$, the set of partially connected clients, together with the sets $F_j$, radius-levels $\ell_j$, inner-balls $B_j$, and the set $D^*$, exactly as defined in Section~\ref{kmedout}; none of these depend on the color of the clients. The only change is in the auxiliary LP: the single coverage constraint~\eqref{const:iter_m} is split into two constraints, one requiring at least $\rho$ red clients to be covered and the other requiring at least $\beta$ blue clients to be covered.

\begin{align} \tag{$\text{LP}_{\text{iter}}'$} \label{LP-iter-col}
\text{Minimize} \quad & \sum_{j \in D_\text{part}} \sum_{i \in F_j} d'(i,j) y_i + \sum_{j \in D_\text{full}} \left( \sum_{i \in B_j} d'(i,j) y_i + (1 - y(B_j)) \Delta_{\ell_j} \right) \\
\text{subject to} \quad & y(F) \leq k, \label{const:iter_col_k} \\
& |D_\text{full} \cap R| + \sum_{j \in D_\text{part} \cap R} y(F_j) \geq \rho, \label{const:iter_col_rho} \\
& |D_\text{full} \cap B| + \sum_{j \in D_\text{part} \cap B} y(F_j) \geq \beta, \label{const:iter_col_beta} \\
& y(F_j) = 1, \hspace{1cm} \forall j \in D^*, \label{const:iter_col_Dstar} \\
& y(B_j) \leq 1, \hspace{0.95cm} \forall j \in D_\text{full}, \label{const:iter_col_B_full} \\
& y(F_j) \leq 1, \hspace{1cm} \forall j \in D_\text{part}, \label{const:iter_col_F_part} \\
& y_i \in [0,1], \hspace{1.08cm} \forall i \in F. \label{const:iter_col_domain}
\end{align}

The vector $y^*$ from Lemma~\ref{lem:eliminate_x_col} is feasible for~\eqref{LP-iter-col} (with $D_\text{part} = D$, $D_\text{full} = \emptyset$, and $D^* = S_0$ initially), and the expected value of the objective function is bounded exactly as in Lemma~\ref{lem:iter_lp_feasible}, i.e., at most $\frac{\tau-1}{\ln \tau}\tilde{U}'$; the argument only uses property (d) of Lemma~\ref{lem:eliminate_x_col} and Lemma~\ref{lem:random_dist}, neither of which involves colors.

\begin{algorithm}[h]
\caption{Iterative Rounding Algorithm (colorful)} \label{alg:iterative_rounding_col}
\begin{algorithmic}
    \Require $I = (F, R, B, d, k, \rho, \beta, S_0)$, $\sigma, \delta \in (0, 1/2)$, $U \ge 0$, $y^* \in [0,1]^F$, $(F_j)_{j \in D \cup S_0}$, and $(\ell_j)_{j \in D \cup S_0}$
    \Ensure a new solution $y^*$ which is almost-integral
\end{algorithmic}
\noindent\rule{\linewidth}{0.4pt}

\begin{algorithmic}[1]
    \State $D_\text{full} \gets \emptyset$, $D_\text{part} \gets D$, $D^* \gets S_0$
    \While{true}
        \State Find an optimum vertex point solution $y^*$ to~\eqref{LP-iter-col}
        \If{there exists some $j \in D_\text{part}$ such that $y^*(F_j) = 1$}
            \State $D_\text{part} \gets D_\text{part} \setminus \{j\}$
            \State $D_\text{full} \gets D_\text{full} \cup \{j\}$
            \State $B_j \gets \{i \in F_j : d'(i,j) \le \Delta_{\ell_j-1}\}$
            \State \Call{Update-$D^*$}{$j$}
        \ElsIf{there exists $j \in D_\text{full}$ such that $y^*(B_j) = 1$}
            \State $\ell_j \gets \ell_j - 1$
            \State $F_j \gets B_j$
            \State $B_j \gets \{i \in F_j : d'(i,j) \le \Delta_{\ell_j-1}\}$
            \State \Call{Update-$D^*$}{$j$}
        \Else
            \State \textbf{break}
        \EndIf
    \EndWhile
    \State \Return $y^*$
\end{algorithmic}
\noindent\rule{\linewidth}{0.4pt}

\begin{algorithmic}[1]
    \Function{Update-$D^*$}{$j$}
        \If{there exists no $j' \in D^*$ with $\ell_{j'} \le \ell_j$ and $F_j \cap F_{j'} \neq \emptyset$}
            \State Remove from $D^*$ all $j'$ such that $F_j \cap F_{j'} \neq \emptyset$
            \State $D^* \gets D^* \cup \{j\}$
        \EndIf
    \EndFunction
\end{algorithmic}
\end{algorithm}

This algorithm is identical to {Algorithm~\ref{alg:iterative_rounding}} of Section~\ref{kmedout}, except that it solves~\eqref{LP-iter-col} instead of~\eqref{LP-iter} at every iteration. In particular, Claim~\ref{claim:iter_invariants} (the basic invariants maintained by the algorithm) holds without any change, since it only concerns the sets $D_\text{full}$, $D_\text{part}$, $D^*$, $F_j$, $B_j$, and $\ell_j$, none of which reference colors.

\subsubsection{Analysis}

Because~\eqref{LP-iter-col} has one extra non-core constraint compared to~\eqref{LP-iter} (two coverage constraints instead of one), the almost-integral solutions it produces may now have up to four strictly fractional variables, instead of two.

\begin{lemma} \label{lem:iter_col_almost_integral}
If after Step 3 in {Algorithm~\ref{alg:iterative_rounding_col}}, neither constraint~\eqref{const:iter_col_B_full} nor~\eqref{const:iter_col_F_part} is tight for $y^*$, then there are at most $4$ strictly fractional $y^*$ variables.
\end{lemma}

\begin{proof}
Assume $y^*$ has $n' \ge 5$ strictly fractional values. We pick $|F|$ independent tight inequalities in~\eqref{LP-iter-col} that define $y^*$. Without loss of generality, we may assume that if $y^*_i \in \{0,1\}$, then the corresponding constraint~\eqref{const:iter_col_domain} for $i$ is picked. Thus, there are exactly $|F|-n'$ tight constraints from among~\eqref{const:iter_col_domain}, and we must pick $n'$ extra tight constraints from among~\eqref{const:iter_col_k}, \eqref{const:iter_col_rho}, \eqref{const:iter_col_beta}, or~\eqref{const:iter_col_Dstar} that are independent from the previous $|F|-n'$.

Each such tight constraint must contain at least two fractional variables. Constraints of the form~\eqref{const:iter_col_Dstar} do not share variables, because for distinct $j,j' \in D^*$ one has $F_j \cap F_{j'} = \emptyset$. Moreover, if constraint~\eqref{const:iter_col_k} is picked, it should be tight and, as it contains all variables, there must be at least two fractional variables in it that are not contained in any picked tight constraint of the form~\eqref{const:iter_col_Dstar}, otherwise the picked constraints would not be independent; the same holds for constraint~\eqref{const:iter_col_rho} with respect to \eqref{const:iter_col_Dstar}. Thus, the number of tight constraints of the form~\eqref{const:iter_col_k} and~\eqref{const:iter_col_rho} together is at most $n'/2$, and the number of tight constraints of the form~\eqref{const:iter_col_k}, \eqref{const:iter_col_rho}, \eqref{const:iter_col_beta}, or~\eqref{const:iter_col_Dstar} is at most $n'/2+2$. This value is strictly less than $n'$ if $n' \ge 5$. Thus, we must have $n' \le 4$.
\end{proof}

The remaining ingredient needed to bound the connection cost of the almost-integral solution — namely, the existence of a nearby unit of open facility for every client of $D_\text{full}$ — does not depend on colors either, and carries over unchanged from Claims~\ref{claim:intersect_distance} and~\ref{claim:Dstar_distance} and Lemma~\ref{lem:full_client_distance}. Together with Lemma~\ref{lem:iter_col_almost_integral}, this yields, exactly as in the proof of Theorem~\ref{thm:solution_iter}, an almost-integral solution $y^*$ (with at most four fractional values) to~\eqref{LP-basic-col} of expected cost at most $\alpha \tilde{U}'$, where $\alpha \approx 7.081$ is the same constant as before.

\subsection{Opening at most $k+1$ facilities}

From the proof of Lemma~\ref{lem:iter_col_almost_integral}, note that if $n' = 4$, then there must be two picked tight constraints among~\eqref{const:iter_col_k} and~\eqref{const:iter_col_rho}, each with two fractional variables, which implies that there are two sets of two fractional variables, each summing to $1$.

Let $y$ be the vector returned by {Algorithm~\ref{alg:iterative_rounding_col}}. We may assume that $y(F) = k$, as otherwise one could open more facilities without increasing the cost. By the observation above, the strictly fractional variables of $y$ can appear only in one of the following four formats:
\begin{enumerate}
    \item Two fractional variables that sum to $1$;
    \item Three fractional variables that sum to $2$;
    \item Three fractional variables that sum to $1$; \label{frac_fac_format_3_col}
    \item Four fractional variables, in two pairs each summing to $1$. \label{frac_fac_format_4_col}
\end{enumerate}

If the fractional values of $y$ appear in formats 1 or 2, we can simply integrally open all the fractionally open facilities, which increases the number of open facilities by exactly one, guaranteeing at most $k+1$ open facilities. For formats~\ref{frac_fac_format_3_col} and~\ref{frac_fac_format_4_col}, we cannot open all the fractional facilities without ending up with $k+2$ open facilities; instead, we choose one fractional facility to close and integrally open the others, which again guarantees at most $k+1$ open facilities. The following two lemmas show how to choose the facility to close so that the coverage requirement for each color still holds.

\begin{lemma}
\label{lem:format3_col}
If the fractional facilities of $y$ appear in format~\ref{frac_fac_format_3_col}, then we can find an integral solution that opens at most $k+1$ facilities and satisfies the coverage requirement for each of the colors.
\end{lemma}
\begin{proof}
Let $i_1,i_2,i_3$ be the three facilities with fractional values of $y$. We have $y_{i_1} + y_{i_2} + y_{i_3} = 1$. For each $X \subseteq \{i_1,i_2,i_3\}$, set
\[R_X = \{j \in D_\text{part} \cap R : F_j = X\} \text{ \quad and \quad } B_X = \{j \in D_\text{part} \cap B : F_j = X\}.\]
By definition of $D_\text{part}$, one has $R_{\{i_1,i_2,i_3\}} = B_{\{i_1,i_2,i_3\}} = \emptyset$, since otherwise, for instance if $j \in R_{\{i_1,i_2,i_3\}}$, constraint~\eqref{const:iter_col_F_part} would be tight.
We will open the facilities $i$ and $i'$ that maximize $|R_{\{i\}}|$ and $|B_{\{i'\}}|$, respectively, and close a remaining one. A client $j$ is served after the rounding if at least one facility of $F_j$ is open.

Suppose, by possibly relabeling the facilities, that $i_3$ is the facility that we will close. This means that $\max\{|R_{\{i_1\}}|,|R_{\{i_2\}}|\} \geq |R_{\{i_3\}}|$ and $\max\{|B_{\{i_1\}}|,|B_{\{i_2\}}|\} \geq |B_{\{i_3\}}|$. Thus, the number of red clients served after the rounding is
\begin{align*}
    |D_\text{full} & \cap R| + |R_{\{i_1,i_2\}}| + |R_{\{i_1,i_3\}}| + |R_{\{i_2,i_3\}}| + |R_{\{i_1\}}| + |R_{\{i_2\}}|\\
    > \ & |D_\text{full} \cap R| + |R_{\{i_1,i_2\}}| + |R_{\{i_1,i_3\}}| + |R_{\{i_2,i_3\}}| + y_{i_1} |R_{\{i_1\}}| + y_{i_3} |R_{\{i_1\}}| + y_{i_2}|R_{\{i_2\}}| + y_{i_3} |R_{\{i_2\}}| \\
    \ge \ & |D_\text{full} \cap R| + |R_{\{i_1,i_2\}}| + |R_{\{i_1,i_3\}}| + |R_{\{i_2,i_3\}}| + y_{i_1} |R_{\{i_1\}}| + y_{i_2}|R_{\{i_2\}}| + y_{i_3} \max \{|R_{\{i_1\}}|,|R_{\{i_2\}}|\} \\
    \ge \ & |D_\text{full} \cap R| + |R_{\{i_1,i_2\}}| + |R_{\{i_1,i_3\}}| + |R_{\{i_2,i_3\}}| + y_{i_1} |R_{\{i_1\}}| + y_{i_2}|R_{\{i_2\}}| + y_{i_3} |R_{\{i_3\}}| \ge \rho,
\end{align*}
where the last inequality holds because of constraint~\eqref{const:iter_col_rho} of~\eqref{LP-iter-col}. Analogously, the number of blue clients served after the rounding is
\begin{align*}
    |D_\text{full} & \cap B| + |B_{\{i_1,i_2\}}| + |B_{\{i_1,i_3\}}| + |B_{\{i_2,i_3\}}| + |B_{\{i_1\}}| + |B_{\{i_2\}}|\\
    \ge \ & |D_\text{full} \cap B| + |B_{\{i_1,i_2\}}| + |B_{\{i_1,i_3\}}| + |B_{\{i_2,i_3\}}| + y_{i_1} |B_{\{i_1\}}| + y_{i_2}|B_{\{i_2\}}| + y_{i_3} \max \{|B_{\{i_1\}}|,|B_{\{i_2\}}|\} \\
    \ge \ & |D_\text{full} \cap B| + |B_{\{i_1,i_2\}}| + |B_{\{i_1,i_3\}}| + |B_{\{i_2,i_3\}}| + y_{i_1} |B_{\{i_1\}}| + y_{i_2}|B_{\{i_2\}}| + y_{i_3} |B_{\{i_3\}}| \ge \beta.
\end{align*}
\end{proof}

\begin{lemma}
\label{lem:format4_col}
If the fractional facilities of $y$ appear in format~\ref{frac_fac_format_4_col}, then we can find an integral solution that opens at most $k+1$ facilities and satisfies the coverage requirement for each of the colors.
\end{lemma}
\begin{proof}
Let $i_1, i_2,i_3,i_4$ be the fractional facilities of $y$. Assume that $y_{i_1} + y_{i_2} = 1 = y_{i_3} + y_{i_4}$. We may assume, after possibly relabeling the variables, that $y_{i_1} = \max\{y_i : i \in \{i_1,i_2,i_3,i_4\}\}$. Thus, $y_{i_2} = \min\{y_i : i \in \{i_1,i_2,i_3,i_4\}\}$. We will never close $i_1$, so, initially, we will choose a facility to close from the set $\{i_{2}, i_{3},i_{4}\}$. For each $X \subseteq \{i_1,i_2,i_3,i_4\}$, set
\[R_X = \{j \in D_\text{part} \cap R : F_j = X\} \text{ \quad and \quad } B_X = \{j \in D_\text{part} \cap B : F_j = X\}.\]
By definition of $D_\text{part}$ and since~\eqref{const:iter_col_F_part} is not tight, one has $R_X = B_X = \emptyset$ for each $X \subseteq \{i_1,i_2,i_3,i_4\}$ with $|X| \geq 3$, and $R_{\{i_1,i_2\}} = R_{\{i_1,i_3\}} = R_{\{i_1,i_4\}} = R_{\{i_3,i_4\}} = B_{\{i_1,i_2\}} = B_{\{i_1,i_3\}} = B_{\{i_1,i_4\}} = B_{\{i_3,i_4\}} = \emptyset$.

From the set $\{i_2,i_3,i_4\}$, we will open the facilities $i$ and $i'$ that maximize $|R_{\{i\}}|$ and $|B_{\{i'\}}|$, respectively, and close a remaining one. A client $j$ is served after the rounding if at least one facility of $F_j$ is open. We focus on the red clients, as the analysis for the blue clients is analogous.

First, suppose $i = i_2$, that is, $|R_{\{i_2\}}| \geq \max\{|R_{\{i_3\}}|,|R_{\{i_4\}}|\}$. We may assume the closed facility is $i_3$, that is, $\max\{|B_{\{i_2\}}|,|B_{\{i_4\}}|\} \geq |B_{\{i_3\}}|$. Then, the number of red clients covered after the rounding is
\begin{align*}
    |D_\text{full} & \cap R| + |R_{\{i_1\}}| + |R_{\{i_4\}}| + |R_{\{i_2,i_3\}}| + |R_{\{i_2,i_4\}}| + |R_{\{i_2\}}| \\
    &=|D_\text{full} \cap R| + |R_{\{i_1\}}| + |R_{\{i_4\}}| + |R_{\{i_2,i_3\}}| + |R_{\{i_2,i_4\}}| + (y_{i_2} + y_{i_1}) |R_{\{i_2\}}| \\
    &=|D_\text{full} \cap R| + |R_{\{i_1\}}| + |R_{\{i_4\}}| + |R_{\{i_2,i_3\}}| + |R_{\{i_2,i_4\}}| + y_{i_2}|R_{\{i_2\}}| + y_{i_1} |R_{\{i_2\}}|\\
    &\geq|D_\text{full} \cap R| + |R_{\{i_1\}}| + |R_{\{i_4\}}| + |R_{\{i_2,i_3\}}| + |R_{\{i_2,i_4\}}| + y_{i_2}|R_{\{i_2\}}| + y_{i_3} |R_{\{i_3\}}| \geq \rho,
\end{align*}
where the first inequality holds because $y_{i_1} \geq y_{i_3}$ and $|R_{\{i_2\}}| \geq |R_{\{i_3\}}|$, and the last inequality holds because of constraint~\eqref{const:iter_col_rho} of~\eqref{LP-iter-col}. Now, suppose that $i \neq i_2$. We may assume $i = i_3$, that is, $|R_{\{i_3\}}| \geq \max\{|R_{\{i_2\}}|,|R_{\{i_4\}}|\}$. Suppose that the closed facility is $i_2$. The number of red clients covered after the rounding is
\begin{align*}
    |D_\text{full} & \cap R| + |R_{\{i_1\}}| + |R_{\{i_4\}}| + |R_{\{i_2,i_3\}}| + |R_{\{i_2,i_4\}}| + |R_{\{i_3\}}| \\
    &=|D_\text{full} \cap R| + |R_{\{i_1\}}| + |R_{\{i_4\}}| + |R_{\{i_2,i_3\}}| + |R_{\{i_2,i_4\}}| + (y_{i_3} + y_{i_4}) |R_{\{i_3\}}| \\
    &=|D_\text{full} \cap R| + |R_{\{i_1\}}| + |R_{\{i_4\}}| + |R_{\{i_2,i_3\}}| + |R_{\{i_2,i_4\}}| + y_{i_3}|R_{\{i_3\}}| + y_{i_4} |R_{\{i_3\}}|\\
    &\geq|D_\text{full} \cap R| + |R_{\{i_1\}}| + |R_{\{i_4\}}| + |R_{\{i_2,i_3\}}| + |R_{\{i_2,i_4\}}| + y_{i_3}|R_{\{i_3\}}| + y_{i_2} |R_{\{i_2\}}| \geq \rho,
\end{align*}
where the first inequality holds because $|R_{\{i_3\}}| \geq |R_{\{i_2\}}|$ and $y_{i_4} \geq y_{i_2}$ by the choice of $i_1$. If the closed facility is $i_4$ instead, we keep the term $y_{i_4}$ in the first inequality and switch the roles of $i_2$ and $i_4$.
\end{proof}

Combining Lemmas~\ref{lem:format3_col} and~\ref{lem:format4_col} with the trivial rounding of formats 1 and 2, we obtain the main integrality result of this section.

\begin{theorem} \label{thm:colorful_main}
Let $y^*$ be an optimum vertex point solution of~\eqref{LP-iter-col} such that neither constraint~\eqref{const:iter_col_F_part} nor~\eqref{const:iter_col_B_full} is tight for $y^*$. Then we can efficiently find an integral solution $y$ that opens at most $k+1$ facilities and covers at least $\rho$ red clients and $\beta$ blue clients.
\end{theorem}

This improves on the $k+3$ open facilities mentioned as a first attempt in Section~\ref{futurework}: choosing which facility to close according to the color-maximizing rule of Lemmas~\ref{lem:format3_col} and~\ref{lem:format4_col}, instead of closing an arbitrary fractional facility, is what allows us to always land on $k+1$.

% TODO: Theorem~\ref{thm:colorful_main} bounds the number of open facilities
% and guarantees coverage, but does not yet bound the *additional connection
% cost* incurred by closing a fractional facility (the colorful analogue of
% Lemma~\ref{lem:convert_cost_partial} in Section~\ref{kmedout}). Check
% whether this is still needed, and if so, add it here.

\subsection{Putting everything together}

% TODO: adapt the combination argument of Section~\ref{kmedout} (guessing $U$,
% choosing $\sigma = \Theta(\eps^{?})$ and $\delta = \Theta(\eps)$, running the
% sparsification of Theorem~\ref{thm:sparsification_col} over all guessed
% instances, and reconnecting the $\rho - \rho'$ missing red clients and
% $\beta - \beta'$ missing blue clients to $S_0$) once the cost bound of the
% previous subsection is available. The final guarantee should be a
% pseudo-approximation: cost at most $\alpha(1+\eps)U$, opening at most $k+1$
% facilities, and covering at least $\rho$ red and $\beta$ blue clients
% (i.e.\ never violating the coverage requirement).
